\documentclass[10pt,a4paper]{article}

% --- Pacchetti ---
\usepackage[utf8]{inputenc}
\usepackage[T1]{fontenc}

\usepackage{amsmath, amssymb}
\usepackage{geometry}
\usepackage{booktabs}
\usepackage{array}
\usepackage{multicol}
\usepackage{enumitem}
\usepackage{xcolor}
\usepackage{mdframed}
\usepackage{titlesec}
\usepackage{hyperref}
\usepackage{fancyhdr}

\pagestyle{fancy}
\fancyhf{}
\fancyfoot[C]{\small\texttt{Maurizio Prazzoli - maurizio.prazzoli1@studenti.unipr.it - Matricola 378685 - https://github.com/maurizioprazzoli/Logica} \quad|\quad \thepage}
\renewcommand{\headrulewidth}{0pt}
\renewcommand{\footrulewidth}{0.4pt}


\geometry{
  top=1.8cm, bottom=1.8cm,
  left=2cm, right=2cm
}

% --- Colori ---
\definecolor{primaryblue}{RGB}{30, 80, 160}
\definecolor{lightblue}{RGB}{220, 232, 250}
\definecolor{warnyellow}{RGB}{255, 248, 220}
\definecolor{warnborder}{RGB}{200, 160, 0}
\definecolor{greenbg}{RGB}{230, 248, 235}
\definecolor{greenborder}{RGB}{40, 150, 60}
\definecolor{redbg}{RGB}{255, 235, 235}
\definecolor{redborder}{RGB}{180, 40, 40}

% --- Stile titoli ---
\titleformat{\section}[block]
  {\large\bfseries\color{primaryblue}}
  {\thesection.}{0.5em}{}
  [\vspace{-0.3em}\rule{\linewidth}{0.8pt}]

\titleformat{\subsection}[block]
  {\normalsize\bfseries\color{primaryblue!80!black}}
  {\thesubsection.}{0.4em}{}

\titlespacing{\section}{0pt}{0.8em}{0.3em}
\titlespacing{\subsection}{0pt}{0.5em}{0.2em}

% --- Ambienti personalizzati ---
\newmdenv[
  backgroundcolor=lightblue,
  linecolor=primaryblue,
  linewidth=1pt,
  topline=true, bottomline=true, leftline=true, rightline=true,
  innertopmargin=4pt, innerbottommargin=4pt,
  innerleftmargin=6pt, innerrightmargin=6pt,
  skipabove=4pt, skipbelow=4pt
]{definizione}

\newmdenv[
  backgroundcolor=warnyellow,
  linecolor=warnborder,
  linewidth=1pt,
  innertopmargin=4pt, innerbottommargin=4pt,
  innerleftmargin=6pt, innerrightmargin=6pt,
  skipabove=4pt, skipbelow=4pt
]{attenzione}

\newmdenv[
  backgroundcolor=greenbg,
  linecolor=greenborder,
  linewidth=1pt,
  innertopmargin=4pt, innerbottommargin=4pt,
  innerleftmargin=6pt, innerrightmargin=6pt,
  skipabove=4pt, skipbelow=4pt
]{esempio}

% --- Comandi matematici shorthand ---
\newcommand{\nec}{\Box}
\newcommand{\pos}{\Diamond}
\newcommand{\entails}{\vDash}
\newcommand{\proves}{\vdash}
\newcommand{\liff}{\leftrightarrow}

% --- Sopprime numerazione sezioni ---
\setcounter{secnumdepth}{2}

% ============================================================
\begin{document}
% ============================================================

\begin{center}
  {\LARGE\bfseries\color{primaryblue} Logica}\\[4pt]
  {\large Logica Proposizionale (PL) e dei Predicati (QL)}\\[2pt]
  {\small !! Le formule possono contenere errori !! - v.1.1\quad|\quad
   \textit{Nota: le regole \textbf{derivate} non sono ammesse all'esame.}}
\end{center}

\vspace{0.5em}
\hrule
\vspace{0.8em}


% ============================================================
\section{Glossario delle Definizioni Teoriche}
% ============================================================

% Ambiente per le voci del glossario
\newcommand{\voce}[3]{%
  \noindent\begin{mdframed}[
    backgroundcolor=lightblue!60,
    linecolor=primaryblue,
    linewidth=0.6pt,
    innertopmargin=5pt, innerbottommargin=5pt,
    innerleftmargin=8pt, innerrightmargin=8pt,
    skipabove=3pt, skipbelow=3pt
  ]
  {\bfseries\color{primaryblue} #1}%
  \ifx&#2&\else\quad{\small\itshape(#2)}\fi\\[2pt]
  #3
  \end{mdframed}%
}

\voce{Forma argomentativa}{}
{Schema astratto ottenuto sostituendo gli enunciati con variabili logiche.\\
\textit{Es.} $P \to Q,\; P \vdash Q$}

\voce{Esemplificazione}{}
{Argomento concreto ottenuto sostituendo le variabili logiche con enunciati reali.\\
\textit{Es.} ``Se piove mi bagno. Piove. Quindi mi bagno.''}

\voce{Forma valida}{}
{Forma argomentativa le cui \emph{tutte} le esemplificazioni sono argomenti validi.}

\voce{Argomento valido}{}
{Argomento di cui \emph{almeno una} forma è valida: è impossibile che le premesse siano vere e la conclusione falsa.\\
\textit{Es.} $P, Q \vdash P$ \quad e \quad $P, Q \vdash R$ (quest'ultimo solo se $R$ segue dalle premesse).\\
\textit{Es.} ``Pm, Ug solo se Pa. Dunque Pm.''}

\voce{Logicamente valido}{}
{La validità dipende \emph{esclusivamente} dal significato delle costanti logiche di PL e QL.
Non dipende dal significato di termini come ``primo'', ``Superman'', ``più alto di''.\\
\textit{Es.} $Fa \vdash \exists x\,Fx$ \qquad \textit{Non es.} $Fa \vdash Q$ (dipende da $Q$)}

\voce{Tautologicamente valido (PL)}{}
{La validità dipende solo dal significato delle costanti logiche di PL.
Non dipende da concetti atomici o predicati.\\
\textit{Es.} $P \land Q \vdash P$\\
\textbf{Gerarchia:} Taut.\ valido $\Rightarrow$ Log.\ valido $\Rightarrow$ Valido}

\voce{Corretto}{\textit{sound}}
{L'argomento è formalmente \textbf{valido} \emph{e} le premesse sono tutte \textbf{vere nel mondo attuale}.\\
\textit{Es.} ``Tutti gli uomini sono mortali. Socrate è un uomo. Dunque Socrate è mortale.''}

\voce{Fbf --- Formula Ben Formata}{}
{Formula sintaticamente corretta, costruita rispettando le regole grammaticali del linguaggio logico.}

\voce{Tautologia}{}
{Fbf vera in \emph{ogni} valutazione possibile delle variabili.\\
\textit{Es.} $P \lor \lnot P$ \quad (tertium non datur)}

\newpage
\voce{Contraddizione}{}
{Fbf falsa in \emph{ogni} valutazione possibile delle variabili.\\
\textit{Es.} $P \land \lnot P$}

\voce{Contingente}{}
{Fbf vera in alcune valutazioni e falsa in altre: il valore dipende dalle variabili.\\
\textit{Es.} $P \to Q$ (valore dipende dai valori di $P$ e $Q$)}

\voce{Tautologicamente equivalenti}{}
{Due formule $\phi$ e $\psi$ tali che $\phi \leftrightarrow \psi$ è una tautologia (stessa tavola di verità).\\
\textit{Es.} $P \to Q \equiv \lnot P \lor Q$}

\voce{Insieme coerente}{}
{Esiste almeno una valutazione che rende \emph{vere contemporaneamente} tutte le formule dell'insieme.\\
\textit{Es.} $\{P,\; P \to Q\}$: basta $P{=}V,\; Q{=}V$}

\voce{Insieme incoerente}{}
{Non esiste \emph{nessuna} valutazione che renda vere tutte le formule insieme.\\
\textit{Es.} $\{P,\; \lnot P\}$: impossibile}

\voce{Funzionalmente completo}{}
{Insieme di connettivi capace di esprimere \emph{qualsiasi} funzione di verità.\\
\textit{Es.} $\{\lnot,\, \land\}$ oppure $\{\text{NAND}\}$}

\voce{Funzionalmente ridondante}{}
{Insieme funzionalmente completo che contiene connettivi superflui (esprimibili dagli altri).\\
\textit{Es.} $\{\lnot,\, \land,\, \lor\}$: il $\lor$ si ottiene da $\lnot$ e $\land$}

\voce{Funzione di verità}{}
{Funzione che a $n$ valori di verità in ingresso associa un unico valore di verità in uscita.
Con $n$ variabili esistono $2^{2^n}$ funzioni di verità distinte.}

\voce{Decidibile}{}
{Esiste un algoritmo (metodo meccanico) che in \emph{tempo finito} risponde Sì/No ad ogni domanda del sistema.\\
\textbf{PL:} decidibile (tavole di verità, alberi semantici).\\
\textbf{QL:} \emph{indecidibile} al primo ordine: se un argomento è \textbf{valido},
la deduzione naturale lo dimostrerà in tempo finito (completezza);
ma se è \textbf{invalido}, la ricerca di un contromodello può non terminare mai.}

\voce{Correttezza}{\textit{soundness}}
{Se una formula è dimostrabile sintatticamente, allora è vera in tutti i modelli.\\
$\text{Se } \Gamma \vdash \phi \;\Longrightarrow\; \Gamma \vDash \phi$\\
Direzione: \textbf{sintassi $\to$ semantica} (dimostrazione $\to$ verità nei modelli)}

\newpage
\voce{Completezza}{\textit{completeness}}
{Se una formula è vera in tutti i modelli, allora è dimostrabile sintatticamente.\\
$\text{Se } \Gamma \vDash \phi \;\Longrightarrow\; \Gamma \vdash \phi$\\
Direzione: \textbf{semantica $\to$ sintassi} (verità nei modelli $\to$ dimostrazione)\\
\textit{Nota:} completezza $\neq$ decidibilità: in QL, se un argomento è \emph{invalido},
non esiste garanzia algoritmica di dimostrarlo.}

\voce{Teorema}{}
{Formula dimostrabile a partire dall'\emph{insieme vuoto} di premesse: $\vdash \phi$.\\
\textit{Es.} $\vdash P \to P$ \quad e \quad $\vdash P \lor \lnot P$}

\voce{Verità logica}{}
{Formula vera in \emph{ogni modello} (PL: ogni valutazione; QL: ogni struttura con qualsiasi dominio).\\
\textit{Es.} $\forall x\,(x{=}x)$: ``tutto è identico a se stesso''}

\voce{Logicamente equivalenti (QL)}{}
{Due formule $\phi$ e $\psi$ tali che $\phi \leftrightarrow \psi$ è una \emph{verità logica}.\\
\textit{Es.} $\lnot\forall x\,Fx \equiv \exists x\,\lnot Fx$}

\voce{Uso}{}
{Si usa l'espressione per parlare del \emph{mondo}.\\
\textit{Es.} Il gatto è sul letto. \quad (si usa ``gatto'' per riferirsi all'animale)}

\voce{Menzione}{}
{Si parla dell'espressione stessa come oggetto linguistico. Richiede le \textbf{virgolette}.\\
\textit{Es.} ``Gatto'' ha cinque lettere. \quad (si menziona la parola, non l'animale)}

\voce{Coreferenziali}{}
{Termini diversi che si riferiscono allo stesso oggetto nel mondo.\\
\textit{Es.} ``Espero'' e ``Fosforo'' (entrambi denotano il pianeta Venere)}

\voce{Salva veritate}{}
{Principio per cui termini coreferenziali sono sostituibili in certi contesti senza alterare il valore di verità.\\
Valido in contesti \textbf{estensionali}; \textbf{non valido} in contesti \emph{intensionali} (modalità, credenze, ecc.).\\
\textit{Es.} ``8 è pari'' (vero). Sostituendo ``il numero dei pianeti'' al posto di ``8'':
``Il numero dei pianeti è pari'' --- valore di verità dipende dal contesto conoscitivo.\\
\textbf{Opacità referenziale:} contesti in cui la sostituzione \emph{non} preserva il valore di verità.}

\voce{Quantifier Shift}{\textit{fallacia}}
{Inversione ingiustificata dell'ordine dei quantificatori: da $\forall x\,\exists y$ si conclude erroneamente $\exists y\,\forall x$.\\
\textit{Premessa valida:} ``Ogni evento ha una causa.'' \; $\forall x\,\exists y\,\text{Causa}(y,x)$\\
\textit{Conclusione fallace:} ``Esiste una causa per tutti gli eventi.'' \; $\exists y\,\forall x\,\text{Causa}(y,x)$}



% ============================================================
\newpage
\section{Struttura degli Argomenti}
% ============================================================

Un argomento è composto da \textbf{premesse} (ragioni) e una \textbf{conclusione}
(l'enunciato derivato).

\subsection*{Indicatori di Premessa}
\textit{Infatti, poiché, perché, giacché, siccome, in quanto, dato che,
visto che, dal momento che, considerato che, la ragione è che.}

\subsection*{Indicatori di Conclusione}
\textit{Quindi, dunque, perciò, pertanto, così, ragion per cui,
di conseguenza, stando così le cose, ne segue che, se ne deduce che,
in conclusione.}

% ============================================================
\newpage
\section{Tipologie di Argomenti}
% ============================================================

\begin{definizione}
\begin{itemize}[leftmargin=*, itemsep=1pt, topsep=0pt]
  \item \textbf{Valido:} Almeno una forma dell'argomento è valida.\\
    Es.\ $P \to Q,\; P \proves Q$
  \item \textbf{Logicamente valido:} Validità dipende \emph{solo} dalle
    costanti logiche di PL/QL.\\
    Es.\ $Fa \proves \exists x\, Fx$
  \item \textbf{Tautologicamente valido (PL):} Validità dipende solo
    dalle costanti logiche di PL.\\
    Es.\ $P \land Q \proves P$
  \item \textbf{Corretto (\textit{Sound}):} Valido \emph{e} premesse
    tutte vere nel mondo attuale.
\end{itemize}
\end{definizione}

\textbf{Gerarchia:}
$\text{Taut.\ valido} \Rightarrow \text{Log.\ valido} \Rightarrow \text{Valido}$

% ============================================================
\newpage
\section{Tipologie di Formule Ben Formate}
% ============================================================

\begin{definizione}
\begin{itemize}[leftmargin=*, itemsep=1pt, topsep=0pt]
  \item \textbf{Tautologia:} Vera in \emph{ogni} valutazione.
    $\entails P \lor \lnot P$
  \item \textbf{Contraddizione:} Falsa in \emph{ogni} valutazione.
    $\entails \lnot(P \land \lnot P)$
  \item \textbf{Contingente:} Vera in alcune valutazioni, falsa in altre.
    Es.\ $P \to Q$
  \item \textbf{Insiemi coerenti:} Esiste almeno una valutazione che
    rende vere \emph{tutte} le formule dell'insieme simultaneamente.
  \item \textbf{Insiemi incoerenti:} Nessuna valutazione rende vere
    tutte le formule.
\end{itemize}
\end{definizione}


% ============================================================
\newpage
\section{Tavole di Verità}
% ============================================================

Le tavole di verità definiscono il comportamento di ciascun connettivo logico
elencando tutti i possibili valori di verità degli operandi e il
corrispondente valore della formula composta.

\subsection{Negazione --- $\lnot P$}

Inverte il valore di verità: vero diventa falso e viceversa.

\begin{center}
\renewcommand{\arraystretch}{1.3}
\begin{tabular}{c|c}
  \toprule
  $P$ & $\lnot P$ \\
  \midrule
  V & F \\
  F & V \\
  \bottomrule
\end{tabular}
\end{center}

\subsection{Congiunzione --- $P \land Q$}

Vera \emph{solo} quando \textbf{entrambi} gli operandi sono veri.

\begin{center}
\renewcommand{\arraystretch}{1.3}
\begin{tabular}{cc|c}
  \toprule
  $P$ & $Q$ & $P \land Q$ \\
  \midrule
  V & V & V \\
  V & F & F \\
  F & V & F \\
  F & F & F \\
  \bottomrule
\end{tabular}
\end{center}

\subsection{Disgiunzione (inclusiva) --- $P \lor Q$}

Vera quando \textbf{almeno uno} degli operandi è vero.

\begin{center}
\renewcommand{\arraystretch}{1.3}
\begin{tabular}{cc|c}
  \toprule
  $P$ & $Q$ & $P \lor Q$ \\
  \midrule
  V & V & V \\
  V & F & V \\
  F & V & V \\
  F & F & F \\
  \bottomrule
\end{tabular}
\end{center}

\subsection{Condizionale (implicazione materiale) --- $P \to Q$}

Falsa \emph{solo} quando l'antecedente è vero e il conseguente è falso.

\begin{center}
\renewcommand{\arraystretch}{1.3}
\begin{tabular}{cc|c}
  \toprule
  $P$ & $Q$ & $P \to Q$ \\
  \midrule
  V & V & V \\
  V & F & F \\
  F & V & V \\
  F & F & V \\
  \bottomrule
\end{tabular}
\end{center}

\begin{attenzione}
Il condizionale materiale è \textbf{vero per vacuità} quando l'antecedente è falso
(\textit{ex falso quodlibet}).
Questa è la riga che crea più confusione: $F \to Q$ è sempre $V$, indipendentemente da $Q$.
\end{attenzione}

\subsection{Bicondizionale --- $P \liff Q$}

Vera quando entrambi gli operandi hanno lo \textbf{stesso} valore di verità.

\begin{center}
\renewcommand{\arraystretch}{1.3}
\begin{tabular}{cc|c}
  \toprule
  $P$ & $Q$ & $P \liff Q$ \\
  \midrule
  V & V & V \\
  V & F & F \\
  F & V & F \\
  F & F & V \\
  \bottomrule
\end{tabular}
\end{center}

\subsection{Riepilogo Compatto}

Tavola riassuntiva di tutti i connettivi binari principali.

\begin{center}
\renewcommand{\arraystretch}{1.3}
\begin{tabular}{cc|cccc}
  \toprule
  $P$ & $Q$ & $P \land Q$ & $P \lor Q$ & $P \to Q$ & $P \liff Q$ \\
  \midrule
  V & V & V & V & V & V \\
  V & F & F & V & F & F \\
  F & V & F & V & V & F \\
  F & F & F & F & V & V \\
  \bottomrule
\end{tabular}
\end{center}


% ============================================================
\newpage
\section{Equivalenze Logiche Fondamentali}
% ============================================================

\subsection{Logica Proposizionale (PL)}

\renewcommand{\arraystretch}{1.3}
\begin{tabular}{@{}ll@{}}
  \toprule
  \textbf{Formula} & \textbf{Equivalente} \\
  \midrule
  $\lnot(P \land Q)$ & $\lnot P \lor \lnot Q$ \hfill \textit{De Morgan 1} \\
  $\lnot(P \lor Q)$ & $\lnot P \land \lnot Q$ \hfill \textit{De Morgan 2} \\
  $P \to Q$ & $\lnot P \lor Q$ \hfill \textit{Def.\ condizionale} \\
  $P \to Q$ & $\lnot Q \to \lnot P$ \hfill \textit{Contrapposizione} \\
  $\lnot(P \to Q)$ & $P \land \lnot Q$ \hfill \textit{Neg.\ condizionale} \\
  $P$ & $\lnot\lnot P$ \hfill \textit{Doppia negazione} \\
  $P \land Q$ & $Q \land P$ \hfill \textit{Commutatività} \\
  $P \land (Q \lor R)$ & $(P \land Q) \lor (P \land R)$ \hfill \textit{Distributiva} \\
  \bottomrule
\end{tabular}

\subsection{Logica dei Predicati (QL)}

\begin{tabular}{@{}ll@{}}
  \toprule
  \textbf{Formula} & \textbf{Equivalente} \\
  \midrule
  $\forall x\, Fx$ & $\lnot\exists x\, \lnot Fx$ \\
  $\exists x\, Fx$ & $\lnot\forall x\, \lnot Fx$ \\
  $\forall x\, \lnot Fx$ & $\lnot\exists x\, Fx$ \\
  $\exists x\, \lnot Fx$ & $\lnot\forall x\, Fx$ \\
  \bottomrule
\end{tabular}

% ============================================================
\newpage
\section{Completezza Funzionale, Decidibilità e Metateoremi}
% ============================================================

\subsection{Completezza Funzionale e Decidibilità}

\begin{itemize}[leftmargin=*, itemsep=1pt, topsep=0pt]
  \item \textbf{Funzionalmente completo:} Insieme di connettivi che
    esprime qualsiasi funzione di verità. Es.\ $\{\lnot, \land\}$
    oppure $\{\text{NAND}\}$.
  \item \textbf{Funzionalmente ridondante:} Completo ma con connettivi
    superflui. Es.\ $\{\lnot, \land, \lor\}$.
  \item \textbf{PL:} Decidibile (tavole di verità, alberi semantici).
  \item \textbf{QL:} \emph{Indecidibile} al primo ordine: argomenti validi
    sono dimostrabili in tempo finito, ma per argomenti invalidi la ricerca
    di un contromodello può non terminare.
\end{itemize}

\subsection{Metateoremi: Correttezza e Completezza}

\begin{definizione}
\begin{itemize}[leftmargin=*, itemsep=2pt, topsep=0pt]
  \item \textbf{Correttezza (\textit{Soundness}):}
    Sintassi $\Rightarrow$ Semantica.
    \[ \text{Se } \Gamma \proves \phi \;\text{ allora }\; \Gamma \entails \phi \]
  \item \textbf{Completezza (\textit{Completeness}):}
    Semantica $\Rightarrow$ Sintassi.
    \[ \text{Se } \Gamma \entails \phi \;\text{ allora }\; \Gamma \proves \phi \]
\end{itemize}
\end{definizione}

Valgono per PL e per QL (dimostrazione di completezza per QL: Gödel, 1930).

\begin{attenzione}
\textbf{Completezza $\neq$ Decidibilità.} Se un argomento è
\emph{invalido}, non esiste un metodo sintattico garantito per
dimostrarlo (si cerca un \emph{contromodello}).
\end{attenzione}

% ============================================================
\newpage
\section{Alberi Semantici}
% ============================================================

Gli alberi semantici (\textit{semantic tableaux}) verificano la validità di un argomento
per \emph{refutazione}: si inserisce nell'albero l'insieme delle \textbf{premesse più la negazione della conclusione}
($\Gamma \cup \{\lnot\phi\}$) e si espandono le formule applicando le regole.

\begin{attenzione}
\begin{itemize}[leftmargin=*, itemsep=2pt, topsep=0pt]
  \item \textbf{Tutti i rami si chiudono} (ogni ramo contiene $\psi$ e $\lnot\psi$ per qualche $\psi$)\\
    $\Rightarrow$ l'insieme $\Gamma \cup \{\lnot\phi\}$ è \textbf{inconsistente}
    $\Rightarrow$ l'argomento è \textbf{valido}.
  \item \textbf{Almeno un ramo resta aperto}\\
    $\Rightarrow$ quel ramo fornisce un \textbf{contromodello} esplicito
    $\Rightarrow$ l'argomento è \textbf{invalido}.
\end{itemize}
\end{attenzione}

\textit{Perché funziona?} Se fosse possibile avere le premesse vere e la conclusione falsa
(cioè $\lnot\phi$ vera), almeno un ramo resterebbe aperto. Chiudere tutti i rami
dimostra che tale scenario è impossibile.

\subsection{Regole NON Ramificanti}

\begin{tabular}{@{}ll@{}}
  \toprule
  \textbf{Formula} & \textbf{Risultato} \\
  \midrule
  $\lnot\lnot P$ & $P$ \\
  $P \land Q$ & $P,\; Q$ \\
  $\lnot(P \lor Q)$ & $\lnot P,\; \lnot Q$ \\
  $\lnot(P \to Q)$ & $P,\; \lnot Q$ \\
  \bottomrule
\end{tabular}

\subsection{Regole Ramificanti}

\begin{tabular}{@{}lll@{}}
  \toprule
  \textbf{Formula} & \textbf{Ramo sx} & \textbf{Ramo dx} \\
  \midrule
  $P \lor Q$ & $P$ & $Q$ \\
  $P \to Q$ & $\lnot P$ & $Q$ \\
  $\lnot(P \land Q)$ & $\lnot P$ & $\lnot Q$ \\
  $P \liff Q$ & $P,\; Q$ & $\lnot P,\; \lnot Q$ \\
  $\lnot(P \liff Q)$ & $P,\; \lnot Q$ & $\lnot P,\; Q$ \\
  \bottomrule
\end{tabular}

\subsection{Regole per i Quantificatori}

\begin{tabular}{@{}lll@{}}
  \toprule
  \textbf{Formula} & \textbf{Risultato} & \textbf{Note} \\
  \midrule
  $\exists x\, P(x)$ & $P(c)$ & $c$ \emph{costante nuova} \\
  $\lnot\forall x\, P(x)$ & $\exists x\, \lnot P(x)$ & riduzione a $\exists$ \\
  $\lnot\exists x\, P(x)$ & $\forall x\, \lnot P(x)$ & riduzione a $\forall$ \\
  $\forall x\, P(x)$ & $P(t)$ & $t$ termine già presente;\\
                     &        & riapplicabile \\
  \bottomrule
\end{tabular}

\subsection{Ordine di Espansione (Priorità)}

\begin{attenzione}
\begin{enumerate}[leftmargin=*, itemsep=1pt, topsep=0pt]
  \item Regole \textbf{non ramificanti} (semplificano senza ramificare).
  \item Regole \textbf{esistenziali} $\exists$ e $\lnot\forall$
    (costante nuova, non ramificano).
  \item Regole \textbf{ramificanti} (moltiplicano i rami).
  \item Regole \textbf{universali} $\forall$ e $\lnot\exists$
    (riapplicabili, usare per ultime).
\end{enumerate}
\end{attenzione}

% ============================================================
\newpage
\section{Deduzione Naturale --- Calcolo Proposizionale}
% ============================================================

\begin{attenzione}
\textbf{Le regole derivate NON sono ammesse all'esame.}
\end{attenzione}

\subsection{Regole Base PL}

% Macro per righe di deduzione: \riga{n}{formula}{giustificazione}
\newcommand{\riga}[3]{$#1$ & $#2$ & \textit{\small #3} \\}
% Macro per righe con ipotesi (barra verticale simulata)
\newcommand{\rigah}[3]{$#1$ & \quad $#2$ & \textit{\small #3} \\}

\renewcommand{\arraystretch}{1.35}

% ------------------------------------------------------------------
\subsubsection*{$\land$I --- Introduzione della Congiunzione}
Da $\phi$ e $\psi$ separatamente, si può scrivere $\phi \land \psi$.

\begin{tabular}{@{}r@{\;\;}l@{\quad}l@{}}
\riga{1}{P}{Pre}
\riga{2}{Q}{Pre}
\riga{3}{P \land Q}{\(\land\)I: 1, 2}
\end{tabular}

% ------------------------------------------------------------------
\subsubsection*{$\land$E --- Eliminazione della Congiunzione}
Da $\phi \land \psi$ si estrae $\phi$ oppure $\psi$.

\begin{tabular}{@{}r@{\;\;}l@{\quad}l@{}}
\riga{1}{P \land Q}{Pre}
\riga{2}{P}{\(\land\)E: 1}
\riga{3}{Q}{\(\land\)E: 1}
\end{tabular}

% ------------------------------------------------------------------
\subsubsection*{$\lor$I --- Introduzione della Disgiunzione}
Da $\phi$ si ottiene $\phi \lor \psi$ (indebolimento).

\begin{tabular}{@{}r@{\;\;}l@{\quad}l@{}}
\riga{1}{P}{Pre}
\riga{2}{P \lor Q}{\(\lor\)I: 1}
\end{tabular}

% ------------------------------------------------------------------
\subsubsection*{$\lor$E --- Eliminazione della Disgiunzione}
Da $\phi \lor \psi$, se da $\phi$ segue $\chi$ e da $\psi$ segue $\chi$, allora $\chi$.

\begin{tabular}{@{}r@{\;\;}l@{\quad}l@{}}
\riga{1}{P \lor Q}{Pre}
\riga{2}{P \to R}{Pre}
\riga{3}{Q \to R}{Pre}
\riga{4}{R}{\(\lor\)E: 1, 2, 3}
\end{tabular}

% ------------------------------------------------------------------
\subsubsection*{$\to$E --- Eliminazione del Condizionale (Modus Ponens)}
Da $\phi \to \psi$ e $\phi$, si ottiene $\psi$.

\begin{tabular}{@{}r@{\;\;}l@{\quad}l@{}}
\riga{1}{P \to Q}{Pre}
\riga{2}{P}{Pre}
\riga{3}{Q}{\(\to\)E: 1, 2}
\end{tabular}

% ------------------------------------------------------------------
\subsubsection*{$\to$I / CP --- Introduzione del Condizionale (Conditional Proof)}
Si apre un'ipotesi $\phi$; se si riesce a derivare $\psi$, si scarica l'ipotesi
e si ottiene $\phi \to \psi$.

\begin{tabular}{@{}r@{\;\;}l@{\quad}l@{}}
\riga{1}{P \to Q}{Pre}
\rigah{2}{\mid\; Q \to R}{Supp}
\rigah{3}{\mid\; R}{\(\to\)E: 1 (\textit{con P})\ldots}
\riga{4}{(Q \to R) \to R}{CP: 2--3}
\end{tabular}

\noindent\textit{Esempio minimo:} $P \vdash Q \to P$

\begin{tabular}{@{}r@{\;\;}l@{\quad}l@{}}
\riga{1}{P}{Pre}
\rigah{2}{\mid\; Q}{Supp}
\rigah{3}{\mid\; P}{Rep: 1}
\riga{4}{Q \to P}{CP: 2--3}
\end{tabular}

% ------------------------------------------------------------------
\subsubsection*{$\lnot$E --- Eliminazione della Negazione (Doppia Negazione)}
Da $\lnot\lnot\phi$ si ottiene $\phi$.

\begin{tabular}{@{}r@{\;\;}l@{\quad}l@{}}
\riga{1}{\lnot\lnot P}{Pre}
\riga{2}{P}{\(\lnot\)E: 1}
\end{tabular}

% ------------------------------------------------------------------
\subsubsection*{$\lnot$I / RAA --- Introduzione della Negazione (Reductio ad Absurdum)}
Si suppone $\phi$; se si deriva $\bot$, si scarica l'ipotesi e si ottiene $\lnot\phi$.

\begin{tabular}{@{}r@{\;\;}l@{\quad}l@{}}
\riga{1}{P \to \lnot P}{Pre}
\rigah{2}{\mid\; P}{Supp}
\rigah{3}{\mid\; \lnot P}{\(\to\)E: 1, 2}
\rigah{4}{\mid\; \bot}{Abs: 2, 3}
\riga{5}{\lnot P}{RAA: 2--4}
\end{tabular}

% ------------------------------------------------------------------
\subsubsection*{Abs --- Assurdo}
Da $\phi$ e $\lnot\phi$ insieme si ricava $\bot$.

\begin{tabular}{@{}r@{\;\;}l@{\quad}l@{}}
\riga{1}{P}{Pre}
\riga{2}{\lnot P}{Pre}
\riga{3}{\bot}{Abs: 1, 2}
\end{tabular}

% ------------------------------------------------------------------
\subsubsection*{EFQ / $\bot$E --- \textit{Ex Falso Quodlibet}}
Da $\bot$ si può derivare qualsiasi formula $\phi$.

\begin{tabular}{@{}r@{\;\;}l@{\quad}l@{}}
\riga{1}{\bot}{Pre}
\riga{2}{Q}{EFQ: 1}
\end{tabular}

% ------------------------------------------------------------------
\subsubsection*{Rep --- Ripetizione}
Una formula già presente nella derivazione può essere ricopiata
(utile all'interno di ipotesi subordinate).

\begin{tabular}{@{}r@{\;\;}l@{\quad}l@{}}
\riga{1}{P}{Pre}
\rigah{2}{\mid\; Q}{Supp}
\rigah{3}{\mid\; P}{Rep: 1}
\riga{4}{Q \to P}{CP: 2--3}
\end{tabular}

% ============================================================
\newpage
\section{Deduzione Naturale --- Calcolo dei Predicati}
% ============================================================

\subsection{Regole per i Quantificatori}

\begin{tabular}{@{}p{2.5cm}p{5cm}@{}}
  \toprule
  \textbf{Regola} & \textbf{Schema e Vincoli} \\
  \midrule
  $\forall$E (Elim.\ univ.) &
    Da $\forall x\, Fx$ si ottiene $Fa$.
    \newline \textit{Nessun vincolo: $a$ può essere qualsiasi.} \\
  $\forall$I (Intr.\ univ.) &
    Da $Fa$ si ottiene $\forall x\, Fx$.
    \newline \textbf{Vincolo:} $a$ deve essere \emph{arbitraria},
    cioè non compare in nessuna assunzione aperta. \\
  $\exists$I (Intr.\ esist.) &
    Da $Fa$ si ottiene $\exists x\, Fx$.
    \newline \textit{Nessun vincolo.} \\
  $\exists$E (Elim.\ esist.) &
    Da $\exists x\, Fx$, con derivazione di $\psi$ dall'ipotesi $Fa$,
    si ottiene $\psi$.
    \newline \textbf{Vincolo:} $a$ deve essere \emph{completamente nuova}:
    non in $\exists x\, Fx$, non in $\psi$, non in assunzioni attive. \\
  \bottomrule
\end{tabular}

\subsection{Regole per l'Identità}

\begin{tabular}{@{}ll@{}}
  \toprule
  \textbf{Regola} & \textbf{Schema} \\
  \midrule
  $=$I (Intr.) & $\proves a = a$ (assioma) \\
  $=$E (Elim.) & Da $a = b$ e $\phi(a)$ si ottiene $\phi(b)$ \\
  \bottomrule
\end{tabular}

% ============================================================
\newpage
\section{Schemi di Traduzione in QL}
% ============================================================

\subsection{Quadrato Aristotelico}

\begin{tabular}{@{}ll@{}}
  \toprule
  \textbf{Enunciato} & \textbf{Formalizzazione} \\
  \midrule
  Tutti i $P$ sono $Q$ &
    $\forall x\,(Px \to Qx)$ \\
  Alcuni $P$ sono $Q$ &
    $\exists x\,(Px \land Qx)$ \\
  Nessun $P$ è $Q$ &
    $\forall x\,(Px \to \lnot Qx)$ \\
  Non tutti i $P$ sono $Q$ &
    $\exists x\,(Px \land \lnot Qx)$ \\
  \bottomrule
\end{tabular}

\subsection{Identità e Numeriche}

\begin{tabular}{@{}p{4cm}p{3.8cm}@{}}
  \toprule
  \textbf{Enunciato} & \textbf{Formalizzazione} \\
  \midrule
  Solo $a$ è $P$ &
    $\forall x\,(Px \liff x{=}a)$ \\
  Solo $a$ e $b$ sono $P$ &
    $\forall x\,(Px \liff (x{=}a \lor x{=}b))$ \\
  Esattamente un $P$ &
    $\exists x\,(Px \land \forall y\,(Py \to y{=}x))$ \\
  Almeno due $P$ &
    $\exists x\,\exists y\,(Px \land Py \land x{\neq}y)$ \\
  Al massimo un $P$ &
    $\forall x\,\forall y\,((Px \land Py) \to x{=}y)$ \\
  Al massimo due $P$ &
    $\forall x\,\forall y\,\forall z\,((Px \land Py \land Pz)$ \\
    & $\to (x{=}y \lor x{=}z \lor y{=}z))$ \\
  Tutti i $P$ tranne $a$ sono $Q$ &
    $Pa \land \lnot Qa \land {}$ \\
    & $\forall x\,((Px \land x{\neq}a) \to Qx)$ \\
  \bottomrule
\end{tabular}

% ============================================================
\newpage
\section{Filosofia del Linguaggio}
% ============================================================

\subsection{Uso vs.\ Menzione}
\begin{itemize}[leftmargin=*, itemsep=1pt, topsep=0pt]
  \item \textbf{Uso:} Si usa l'espressione per parlare del mondo.
    Es.\ \textit{Il gatto è sul letto.}
  \item \textbf{Menzione:} Si parla dell'espressione stessa
    (richiede le virgolette).
    Es.\ \textit{``Gatto''} ha cinque lettere.
\end{itemize}

\subsection{Coreferenzialità e Salva Veritate}
\begin{itemize}[leftmargin=*, itemsep=1pt, topsep=0pt]
  \item \textbf{Coreferenziali:} Termini diversi che denotano lo stesso
    oggetto. Es.\ ``Espero'' e ``Fosforo'' (entrambi = Venere).
  \item \textbf{Salva veritate:} Termini coreferenziali sono
    sostituibili in contesti \emph{estensionali} senza cambiare il
    valore di verità. \textbf{Falso} in contesti intensionali
    (modalità, credenze, ecc.).
\end{itemize}

\subsection{Fallacia del Quantifier Shift}

\begin{attenzione}
Inversione ingiustificata dei quantificatori.\\
\textit{Premessa:} $\forall x\,\exists y\,(\text{y causa x})$\\
\textit{Conclusione fallace:} $\exists y\,\forall x\,(\text{y causa x})$
\end{attenzione}


% ============================================================
\newpage
\section{Riepilogo: Semantica dei Modelli (QL)}
% ============================================================

\begin{definizione}
\textbf{Modello} $\mathcal{M} = \langle D, I \rangle$:
\begin{itemize}[leftmargin=*, itemsep=1pt, topsep=0pt]
  \item $D$: dominio non vuoto di individui.
  \item $I$: funzione di interpretazione che assegna denotata a
    costanti individuali e estensioni a lettere predicative.
\end{itemize}
\textbf{Verità logica:} formula vera in \emph{ogni} modello.\\
\textbf{Equivalenza logica (QL):} $\phi \liff \psi$ è una verità logica.\\
\textbf{Teorema (QL):} formula dimostrabile dall'insieme vuoto di premesse.
\end{definizione}



\vspace{0.5em}
\hrule
\vspace{0.3em}

\end{document}
